%!TEX root = ../../thesis.tex

\section{What the Approach Demonstrates and What It Costs}\label{sec:discussion-interpretation}

The agent learns to fuzz, and it does not fuzz competitively.
Both halves of that sentence rest on \cref{chap:results}, and the second is not a caveat attached
to the first but a result of the same experiments.
This section states what the evidence carries, with the target, the baseline and the budget named
for each claim, and then what the approach costs to obtain it.

Credit assignment across a sequence of decisions is the capability the input-action targets
isolate, and the learner has it.
On high-and-low the agent plays high-coverage bytes at a rate above \num{0.99} where a random
policy reaches $\frac{4}{256}\approx\num{0.0156}$~(\cref{sec:results-high-and-low}).
On sequence it recovers the entire 16-byte passcode one position at a time, and the cost of each
additional byte grows from three positions within the first six iterations to 55 iterations for
the last one alone~(\cref{sec:results-sequence}).
That last number is the point: a mutation-based fuzzer with no notion of return has no mechanism
that makes the sixteenth byte cheaper than exhaustive search over it, and the agent's is a value
function that carries reward backwards across steps.
The scope of the claim is narrow.
Both targets reward the agent for progress that no coverage signal on a real program would supply
unprompted (\cref{sec:target-sequence}), and each ran on a single PRG seed.
What is demonstrated is credit assignment under a designed reward, not fuzzing.

On the corpus-action targets the reward is raw coverage, and the agent still converges on
structure.
The highest-coverage cJSON input is \code{[[[[[2,[[[[[1\ldots}\ and the highest-coverage libxml2
input nests tags to a comparable depth~(\cref{tab:best-inputs}), neither of which was described to
the agent in any form.
The structural byte rate makes the same point over training rather than at its end: it rises from
\num{0.08} to a peak of \num{0.90} on cJSON and from \num{0.05} to \num{0.75} on
libxml2~(\cref{fig:structural-decay}).
A recursive-descent parser rewards nesting with coverage, and the agent finds that out from the
coverage vector alone.

Planning also beats undirected mutation at an identical execution budget, on the one target where
that comparison was run.
At \num{1638400} guest executions the cJSON baseline reaches a median of \num{645} covered blocks
against \num{542} for the structure-blind random control, and the campaign's registered primary
metric --- the fraction of episodes whose best seed reaches the 99th percentile of the pooled
control distribution --- separates them by a factor of roughly \num{71}, \SI{76.7}{\percent}
against \SI{1.08}{\percent}, on every paired seed~(\cref{sec:target-cjson}).
The best-performing arm reaches \num{958} blocks, so the gap to the control widens further under a
better configuration~(\cref{tab:campaign-inventory}).
Two things bound this.
The control is uniform random mutation, not a coverage-guided fuzzer, so the comparison establishes
that the learned policy is doing something rather than that it is doing something better than the
state of the art.
And the primary metric's threshold sits below the plateau every trained arm reaches early, so it
separates trained from untrained decisively and the trained arms from each other not at all.

Where coverage does not grow with input structure, the approach buys nothing.
picohttpparser ends at \num{81} covered blocks with a seed-to-seed range of one block, reaches
\SI{90}{\percent} of that ceiling within \SI{25}{\percent} of its budget and gains nothing at all in
the final quarter~(\cref{fig:saturation,tab:campaign-inventory}).
This is not a failure of learning.
The structural byte rate on picohttpparser rises from \num{0.03} to a peak of
\num{0.73}~(\cref{fig:structural-decay}), so the agent learns to produce better-formed requests
just as it does on the other two targets; the target simply does not pay for them, because a
header parser's block count is bounded by its grammar rather than by the depth of any single
request.
The property that predicts whether the approach helps is therefore a property of the target and
not of the agent: coverage that keeps growing as inputs grow more structured.
A practitioner can check it cheaply, and should, before spending anything on the machinery this
thesis describes.

The dominant cost is throughput.
Under the eight-way concurrency the campaigns ran at, a cJSON run executes \num{1638400} inputs in
\SI{21.3}{\hour} and the two 400-iteration campaigns execute \num{3328128} in \SI{42.1}{} and
\SI{40.7}{\hour}, which is between \num{21} and \num{23} guest executions per second in every
case~(\cref{tab:campaign-inventory}).
Removing the search entirely raises this only to \num{41}: the random control covers the same
budget in \SI{11.3}{\hour}, so emulation and the learner, not the planner, set the ceiling.
Measured alone on the card a run is roughly three times faster, which moves the figure to the order
of \num{100} and no further.
A native persistent-mode harness on a parser of this size executes on the order of \num{e4} inputs
per second on one core~\cite{fioraldi_afl_nodate}, a figure this thesis cites rather than measures,
since no conventional fuzzer was run on these targets.
Taken at that value the agent must extract roughly three orders of magnitude more coverage per
execution than random mutation to break even in wall-clock, and the factor of \num{71} it does
achieve on cJSON is measured per execution against a control that shares its own slow harness.
\todo{Pin the \num{e4} figure to a specific measurement in the cited source, or replace it with a
measured AFL++ run on cJSON.}

The search is the component that pays for itself least.
Coverage is flat across simulation counts --- \num{687}, \num{640} and \num{645} blocks at 50, 100
and 200 simulations --- while wall-clock scales close to linearly with them, at \SI{7.6}{},
\SI{11.5}{} and \SI{21.3}{\hour}~(\cref{tab:campaign-inventory}).
In the execution currency the three arms are indistinguishable; in the wall-clock currency the
cheapest of them dominates the configuration this thesis otherwise treats as its baseline.
The policy-improvement diagnostic points the same way: the mean divergence between the prior and
the \gls{mcts} posterior falls from \num{0.99} to \num{0.28} nats over
training~(\cref{fig:policy-improvement}), so the search corrects a policy that increasingly does
not need correcting.
None of this shows that planning is worthless, because the smallest search that was run is still 50
simulations wide; it shows that beyond that width the thesis paid for search and got coverage that
was already there.

Reliability is the second cost, and it is a property of every run rather than of unlucky ones.
The structural byte rate peaks before the end of training on all three corpus targets and finishes
below its peak in each: \num{0.90} to \num{0.80} on cJSON, \num{0.75} to \num{0.55} on libxml2 and
\num{0.73} to \num{0.36} on picohttpparser~(\cref{fig:structural-decay}).
Per-position action entropy falls to between \num{0.4} and \num{0.7} nats by the end of
training~(\cref{fig:collapse-entropy}), and the share of replay transitions carrying a non-zero
reward falls as the reachable coverage is exhausted~(\cref{fig:buffer-reward-density}).
A fuzzer that is best some fraction of the way into its run and worse afterwards needs a stopping
rule that this thesis does not provide, and the comparison it loses is against a technique whose
failure mode is to make no further progress rather than to give up ground.

The remaining costs are ones the design fixed in advance.
The strongest results in this thesis are on the two targets whose reward was engineered, and that
effort transfers to a new target no better than the per-target harness work that coverage-guided
fuzzing was invented to avoid~(\cref{sec:fuzzing-limitations}).
The action space is a flat categorical over an enumerated branching factor with no
masking~(\cref{sec:mdp-actions}), which is what rules out variable-length inputs and
phase-dependent legal actions; the observation is a fixed hashed coverage vector that aliases
blocks and does not grow with the target~(\cref{sec:impl-coverage}); and the corpus holds a single
seed~(\cref{sec:impl-corpus}), which makes selection and eviction degenerate.
Against a fuzzer that ships as one binary, the approach needs an emulator, a learner, a probe stack
and a GPU per run.

Separating the two kinds of limitation is what makes the rest of this useful.
Throughput under emulation, the single-seed corpus, the fixed input length, the hashed observation
and the missing stopping rule are limitations of this implementation, and each has a concrete
remedy that costs engineering rather than insight~(\cref{chap:future-work}).
The dependence on a target whose coverage grows with input structure, the dependence on a reward
that someone designed, and the cost of a learned policy that must be paid before it returns
anything are limitations of the approach.
Only the latter three decide whether reinforcement learning belongs in a fuzzer at all.
